\section{Purpose, checkpoint, and contribution}
\label{sec:scope}

The aim of this continuation is to improve operations that an eventual
efficient geometric recognizer must execute, while making their mathematical
and implementation contracts precise. The result has two parts. The first
accelerates an existing independently checked operation. The second supplies
a missing geometric input to an earlier graph algorithm. Their relationship
is direct: the new geometric observer is itself built from weighted orbit
queries, so it benefits from improved weighted replay without weakening the
source checks.

\subsection{The present complexity target}

Lackenby's February 2021 announcement states an
\(n^{O(\log n)}\) unknot recognition bound in the crossing number
\cite{announcement,talk}. His July 2026 preprint gives a hierarchy algorithm
for incompressibility and unknot recognition. Section~9 explicitly leaves
some operation costs insufficiently specified for a time estimate.
Proposition~9.1 gives the iteration bound \(L(g+1)^L\);
Proposition~10.2 bounds refined hierarchy length by
\(4c_q(\mathcal H)\). Theorems~14.1 and 14.4 provide efficient hierarchy
encoding and compressed cutting results \cite{lackenby26}. These are
substantial existing results. This paper neither reconstructs a complete
proof of the announced bound nor attributes a new polynomial weighted-orbit
theorem to itself.

Classical Agol--Hass--Thurston orbit counting and its weighted extension are
the mathematical foundation of the local operations studied here
\cite[Theorems~12 and 16]{aht}. The question addressed by our first part is
more specific: how much of the current weight partition must an independently
checked replay actually revisit?

\subsection{Encoding and notation}

The global target uses the diagram crossing number \(n\). For a local
normal-surface query, \(T\) is the tetrahedron count, \(B\) the maximum
normal-coordinate bit length, and \(k\) the number of explicitly labelled
components. For an orbit query, \(N\) is the numerical interval size,
\(R\) the number of canonical constant runs, and \(d\) the dimension of a
weight vector. Endpoint costs depend on \(\log(N+1)\), not on \(N\)
literal iterations. A certificate is itself input to its checker, so its
encoded length is charged whenever an arbitrary supplied trace is
considered. These distinctions prevent an operation count in a large
normal coordinate from being mistaken for a polynomial encoded bound.

\subsection{What was inspected}

The implementation checkpoint is \BaseCommit, the \code{main} revision
observed at the beginning of this work. All patches and reproduced
measurements in the package refer to that fixed source. Later repository
changes are not silently incorporated into the baseline.

The inspected source includes the native \code{fastunknot} modules for
weighted histograms, weighted representative transversals, normal surface
geometry, raw cut cells, cut complement components, parallelity bundles,
ordered boundary interfaces, and supplied-cut handlebody certification. The
maintained source and reports distinguish these supplied geometric queries
from the complete diagram recognizer \cite{repo}.

The incoming directory at this checkpoint contains four relevant packages:
\emph{Exceptional Rows}, \emph{Matroidal Dual Forests},
\emph{Nonlinear Activation Forests}, and \emph{Ordered Interface Kernels}.
Their source and integration descriptions were inspected for overlap.
They address finite forest/cancellation selection and related replay
optimizations. None supplies the lazy receiver operation or the native
component-labelled chamber exporter delivered here. Earlier reports do,
however, prove several tempting candidate results already. These
are used with attribution.

\begin{table}[htbp]
\centering\small
\begin{tabular}{L{0.24\textwidth}L{0.34\textwidth}L{0.32\textwidth}}
\toprule
Topic & Earlier result & Contribution of this continuation\\
\midrule
Weighted orbit counting &
Polynomial encoded computation, classical AHT &
A faster replay of the same checked traces\\
\addlinespace
Native fold run growth &
\(R_T\le R_0+2T\), already in \code{suffix_folds.tex} &
Use this bound to control a lazy balanced representation\\
\addlinespace
Receiver-local arithmetic &
Flat producer work \(O(R)+O((S+P)d)\) &
Remove the receiver-run term \(P\) from tree and vector updates\\
\addlinespace
Chamber graphs &
One-census compiler and cographic nonseparation in report~92 &
Authenticate labels and side locations from actual normal coordinates\\
\addlinespace
Hierarchy resets &
Weighted dependency-path bounds in reports~23 and~92 &
No replacement theorem asserted; identify the still-missing geometric bridge\\
\bottomrule
\end{tabular}
\caption{Attribution and scope. \(R\) is the total number of weight runs,
\(S\) the deleted suffix runs, \(P\) the receiver runs, and \(d\) the vector
dimension.}
\label{tab:scope}
\end{table}

\subsection{The first main result}

The earlier suffix-fold optimization already avoids some vector additions
outside the receiving interval, but its state is a flat sequence
\cite{suffix}. A receiver containing many different vectors is still
visited throughout. The checker reconstructs even more of the flat
partition. This is a representation cost, not a necessary part of checking
the topological source.

Our key observation is elementary but useful: adding the same vector to all
runs of an interval preserves every equality and inequality between
neighbouring runs inside that interval. An exact checker can therefore
store that addition at the root of a balanced subtree. It need only split
at the interval endpoints and inspect the two resulting seams.

Theorem~\ref{thm:lazy-fold} makes this observation quantitative. In
particular, the work does not grow with the number of pre-existing receiver
runs. The complete source family in Theorem~\ref{thm:star} shows why this
changes an asymptotic term rather than only a constant factor. The
implementation preserves the residue-anchor and discrete-derivative
checker, so production and replay still use different transport arithmetic.

\subsection{The second main result}

A finite graph can encode which components of a disjoint two-sided surface
system may be cut simultaneously without disconnecting the ambient manifold.
The graph alone is easy to use. The difficult interface question is whether
its vertices and edges genuinely describe the supplied source.

The native exporter takes one embedded normal surface together with an
explicit representative disc for each of its claimed connected components.
It verifies that the representatives are distinct and exhaustive
\emph{as component labels}, verifies two-sidedness, and locates both sides
of each component in the actual cut complement. A successful certificate
therefore supplies the missing geometric contract of report~92 for this
input class \cite{persistence}. Its checker derives side locations through
a separate raw cut-cell construction rather than importing the exporter.

The number of components \(k\) is an explicit input/output parameter. A
normal vector can encode an enormous number of parallel components. This
exporter does not compress an exponentially large arbitrary labelled graph
into a constant-size answer. For a connected large-coordinate surface,
however, \(k=1\), and a single representative suffices regardless of the
number of normal discs.

\subsection{What the results establish}

The local replay theorem is unconditional for the specified checked
interval trace and exact weight representation. The native chamber theorem
is unconditional under the implemented geometric source contract:
a validated finite connected compact orientable triangulation with one
torus boundary, an admissible normal vector, and correctly authenticated
component representatives whose two-sidedness gate succeeds.

The chamber theorem does not require the components to be incompressible,
and does not prove that they are. The resulting graph certifies
nonseparation only. Neither a connected complement nor a graph loop is a
handlebody certificate. Moreover, the graph does not yet contain a
geometric path crossing a particular component once, or a support
certificate showing that such a path survives a later edit.

For the same reason, the new replay engine does not alter the correctness
criterion of the complete recognizer. It reduces work needed to check
existing local certificates. A performance claim about a supplied normal
cut is reported as such, rather than as a speed-up for arbitrary knot
recognition.

\subsection{Novelty in the appropriate sense}

Balanced sequence trees, delayed range additions, normal-surface
connectivity, and dual graphs are established ideas. No priority claim is
made for those foundations. The contributions here are the exact
receiver-independent replay construction with its native trace separation,
the source-derived chamber-label interface and its independently reconstructed
checker, and the executable evidence connecting them to this implementation.
An exhaustive literature priority search for these combinations has not
been performed. The proofs should be judged on their stated claims, and the
measurements on their executed scopes.
